% =====================================================================
% Replacement prose for the integrated manuscript (Fable 5.2 cross-review)
% =====================================================================
% This file is \input-able. It needs the shared transcript-card macros, the
% incident macros of full_drafts/fable_5_2/preamble.tex, and three quote
% files: sections/transcript_quotes.tex (shared, 28 fields),
% full_drafts/opus/sections/opus_quotes.tex (G2, G11) and
% full_drafts/fable_5_2/sections/quotes_supplement.tex (G15, G1, G11, G16),
% plus reviews/fable_5_2/review_quotes.tex (two inline excerpts, G14 and G16,
% generated by tools/build_review_quotes.py).
% reviews/fable_5_2/build/replacement_build.tex compiles it stand-alone.
%
% It replaces only material that changes. Each block names where it goes and
% which draft supplies the unchanged text around it (integration_proposal.md
% gives the full keep/cut list). No transcript text is typed here: every
% excerpt is a generated macro. The one typed quotation is the prompt
% sentence in R2, checked against a saved standard system prompt.
%
% Provenance of numbers. [R] recomputed for this review from the retained
% ledgers or canonical decision records (reviews/fable_5_2/tools, outputs in
% reviews/fable_5_2/data). [I] inherited from the 29 October audit and not
% re-executed: every engine result (payment witnesses, reordered build plans,
% the G7 search, replay) and every semantic code. Appendix block R12 prints
% the same split for readers.
% =====================================================================

% ---------------------------------------------------------------------
% R0. TITLE AND ABSTRACT (replace all four)
% ---------------------------------------------------------------------
\section*{Proposed title and abstract}

\noindent\textbf{Bargaining, building and bankruptcy: four configured LLM systems in twenty complete games of Monopoly}

\medskip
\noindent An agent that trades against rivals has to judge what a deal gives the other side, turn what it acquires into income, and stay solvent. We follow these steps through twenty complete four-player Monopoly games played by four configured LLM systems, using the engine's legal menus, every response attempt, the prompts actually delivered and the recorded events. The four systems complete 17, 18, 15 and 13 color groups by trade and build on 16, 17, 7 and 3 of them at the first opportunity. Groups built at the first opportunity later collect \$53,810 of the \$57,378 in rent paid on trade-completed groups. The twenty groups left unbuilt divide into nine whose first management window never listed construction, four engine defaults, four responses whose text and tool call disagree, and three end-of-turn choices with construction listed. All 61 defaulted build requests of one system were executable after legal reordering. In bargaining, an offer that would complete only its author's color group is accepted 19 times in 289, and one of 187 such opening offers is accepted by a response whose text and call agree. The contrast with offers that complete the recipient's group halves once position in the thread is held fixed, and it is absent for one system. Fifty-six of sixty eliminations could not have been paid from the debtor's own assets, three could, and one was blocked by bank stock. The eventual winner is the creditor in 35 of the 40 eliminations before the last. Seven incidents with verbatim excerpts show the systems bargaining over a rival's building cash, solvency and estate. The study is observational. It estimates no effects and ranks no models, and two of the four systems ran under response conditions that limit comparison.

% ---------------------------------------------------------------------
% R1. INTRODUCTION (complete replacement)
% Base: Fable paragraphs 1-2 for the opening, Opus for the G2 vignette.
% ---------------------------------------------------------------------
\section{Introduction}
\label{sec:intro}

An agent that trades against rivals faces three problems in sequence. It must judge a deal by what the deal gives the other side as well as by what it gets. It must turn what it acquires into income, which takes cash, a legal opportunity and a correctly specified action. And it must stay solvent, because a debt it cannot pay hands what it built to someone else. Evaluations of language-model agents usually examine one of these problems at a time, in a bargaining game, a business simulation or a planning task.

Monopoly puts all three on one path. A player completes a color group, often by trade. Houses multiply the group's rent. A rival who lands there and cannot pay goes bankrupt, and in the implementation studied here the creditor inherits the debtor's deeds with their buildings intact. We ask what four LLM systems did and said at each step of that path in twenty complete games.

One thread from game 2 shows the path. At turn 45 Sol holds the deed Opus needs to complete the yellow group. Opus offers \$325. Sol asks \$650, writing that the lower price does not compensate for handing over the monopoly, and refuses the \$450 that Opus calls final. Twenty-four turns later Opus is down to \$17 and sells Sol its two yellow deeds for two railroads and \$100. Opus is eliminated three turns after that. Sol builds two houses, collects \$44 on the group, and is eliminated by rent owed to Gemini. Gemini inherits the yellow group, builds six houses on it and wins (Incident~\ref{inc:g2}).

The data are twenty complete four-player games among four configured systems, which we call Sol, Gemini, Opus and Grok after their underlying models. The games contain 6,989 decisions and all end by bankruptcy without a turn limit. For every decision the record holds the legal menu, each response attempt, the applied action and the resulting events. We reconstruct all 526 negotiation threads, all 63 color groups completed by trade, all 60 eliminations and every transfer of cash and deeds. Four results organize the paper.

\emph{Execution decides what a negotiated group is worth} (Section~\ref{sec:development}). The four systems complete similar numbers of groups by trade. Two build on nearly every one at the first opportunity and two build on fewer than half, and the groups built at the first opportunity collect 94 percent of the rent later paid on trade-completed groups. The unbuilt groups have four different histories: a stated reason to hold cash, too little cash to act, a build plan that the engine rejected and replaced with an end of turn, and a tool call that ended the turn while its text asked to build.

\emph{Offers that complete the author's group are seldom accepted, and the pooled contrast overstates how distinctive that is} (Section~\ref{sec:bargaining}). Such offers are accepted 19 times in 289. Offers that complete the recipient's group are accepted 24 times in 140 and are mostly countered. About half of the difference in rejection reflects where the two kinds of offer fall in a thread, three of the four systems show it, and counteroffers reprice a requested deed in the same way whether or not it completes a group. Two systems state reasons about the rival's position when they refuse.

\emph{Almost no system gives up with a way to pay} (Section~\ref{sec:liquidity}). Fifty-six of sixty eliminations fall below a generous bound on what the debtor's own assets could raise. Two of the three payable debts were abandoned by a system whose note left out the mortgages that open up once buildings are sold.

\emph{Estates follow rent} (Section~\ref{sec:inheritance}). The eventual winner is the creditor in 35 of the 40 eliminations before the last, and the groups those estates complete collect little rent of their own. Inheritance in these games consolidates a lead that developed groups had already produced. The systems still bargain over it, as when a seller prices a deed as insurance against a rival's inheritance and a likely creditor withdraws an offer from a nearly bankrupt seller.

These results matter for anyone who expects language-model agents to hold and trade assets. A count of deals closed and a count of games won would both misdescribe these four systems, because the value of a deal was set afterwards, at the execution layer and in the cash position the deal left behind. The cohort also shows how much of an outcome can come from the dice. In seven of the twenty games the winner is the player whose seat bought an early monopoly from the bank, whichever system sat there (Section~\ref{sec:cohort}).

This is an observational study of four complete configurations, including their prompts, memory, retry policies and provider behavior. Two configurations ran under conditions that change what their records mean, and Section~\ref{sec:setup} states those conditions before any result. The counts are censuses of a finite, dependent cohort. We do not rank models, estimate the effect of a rival's completion on refusal, or claim that any trade, purchase or bankruptcy was the best available play. Section~\ref{sec:discussion} maps the evidence onto eight broader questions about strategic agents and says which ones this cohort cannot answer.

The work builds on economic and bargaining evaluations that record process as well as outcome~\cite{backlund2025vendingbench,li2026arena,wu2026agentmonopoly,hu2026finboardbench,bianchi2024negotiationarena}. It adds a connected account of one complete economy, from every offer through every traded monopoly to every bankruptcy.

% ---------------------------------------------------------------------
% R2. SETUP (replacement paragraphs)
% Base: Opus sections/setup.tex (rules list, configuration table, conditions
% table). Replace the refusal row text, the honesty sentence and the cohort
% paragraph with the three blocks below.
% ---------------------------------------------------------------------
\section{Environment, cohort and conditions (replacement paragraphs)}
\label{sec:setup}

\paragraph{Refusal and recovery.} Every decision has a budget of two response attempts. After an invalid action the second attempt is a corrective request that quotes the validator's message. For Opus the provider refused 1,605 attempts in 1,605 of its 1,696 decisions, returning the same Terms-of-Service message each time. In 1,604 of these the refusal is the first attempt, and the second attempt is an action-only request that drops the private-note field and the sentence asking for honest private notes. That request spends the budget, so an invalid action-only response is followed by the engine's fallback and not by a correction. In the remaining decision, G2 d129, an invalid build plan comes first and the refusal second. The refusal affects 215 of 306 Opus decisions in games 1 to 4 and all 1,390 in games 5 to 20. We do not know what triggers it. The change at game 5 coincides with a source revision that touched no engine, prompt, resolver or client file, and the record does not identify a cause.

\paragraph{What the prompt asks of each channel.} The standard prompt states that ``Deception in public chat is allowed; private thoughts should reflect true reasoning.'' A false public statement therefore breaks no instruction. A false private note does not by itself show dishonesty or a breach of the request, because the request concerns the reasoning and a mistaken belief can be honestly reported. The action-only request removes the private note and the sentence about it.

\paragraph{Cohort and dependence.}\label{sec:cohort} Five seeds are each played in four cyclic seat rotations, so every system sits in every seat five times. Games run from 75 to 386 turns. Sol wins eight, Gemini seven, Opus three and Grok two.

These totals are not a ranking, and the cohort shows why. Rotations within a block share their early dice. In block 3 the player in seat 2 buys both dark-blue deeds from the bank by turn 29 in all four rotations. Each of the four systems sits in that seat once, and each wins from it. In block 4 the player in seat 1 completes the orange group by purchase on turn 47 in three of the four rotations, under three different systems, and wins all three. In block 2, by contrast, Sol wins from every seat. Seven games thus follow the seat that drew an early monopoly and four follow one system across seats. Twenty games in five blocks cannot separate these influences. We attach no confidence intervals to dependent decisions and draw no seat effect from two blocks.

% ---------------------------------------------------------------------
% R3. MEASUREMENT (replacement paragraphs)
% Base: Sol "Units, reconstruction and economic endpoints" and "Checked
% evidence and inherited validation"; Opus "Legal-sequence tests".
% ---------------------------------------------------------------------
\section{Measurement (replacement paragraphs)}
\label{sec:measurement}

\paragraph{Offers.} The offer ledger has one row for each of the 920 times a player had to answer a trade offer. The \emph{author} is the player who wrote the terms on the table, who after a counter is the former recipient. We classify each offer by who would newly complete a color group if it were accepted: the author only, the recipient only, both or neither. Completion means ownership of every deed in a color group, mortgaged or not. It does not mean the group can be built on. Acceptance is on the legal menu for 918 offers, and rates that require a real choice to accept use those. Three further descriptions of an offer turn out to matter. Its \emph{position} is whether it opens a thread or follows a counter. Its \emph{deed direction} is which side would receive more deeds. A counter is a \emph{pure repricing} when it keeps every deed and card of the offer it answers and changes only the cash.

\paragraph{Reading text.} Two results rest on coding by language models. Two readers coded the stated reasons in the 74 outright rejections of offers that would complete the recipient's group, and one reader coded forward-looking statements in a fixed sample of 240 decisions. The second rejection reader had seen aggregate counts and one row's labels. No person annotated either set, and agreement between model readers does not validate a code. Grok and Gemini wrote 62 of the 74 rejections, so the coded reasons describe mainly those two systems.

% ---------------------------------------------------------------------
% R4. BARGAINING (complete replacement of the section)
% ---------------------------------------------------------------------
\section{Bargaining over completion}
\label{sec:bargaining}

\subsection{Volume, roles and settlement}

The twenty games contain 526 negotiation threads and 920 responses. One thread in five settles (113). The median thread has one response and the longest has eleven counters, no thread spans a turn, and 414 threads follow an earlier thread between the same two players. In 44 settlements the accepted terms were written by the player who did not open the thread.

The four systems take different roles. Opus opens 240 threads, Grok 136, Sol 132 and Gemini 18, and 44, 39, 19 and 11 of these settle. Grok's count carries a condition from Section~\ref{sec:setup}. Sixty-four of its 136 threads were opened by a tool call that proposed a trade while the accompanying text named another action, usually ending the turn. Seventeen of those threads settled. Recipients answered these offers like any others, so they stay in the offer population, but Grok's count of openings is not a clean measure of decisions to negotiate.

\subsection{Responses by who would complete a group}
\label{sec:cells}

\begin{table}[t]
\centering\small
\setlength{\tabcolsep}{4.2pt}
\begin{tabular}{@{}lrrrrrrrrrrrr@{}}
\toprule
 & \multicolumn{4}{c}{All offers} & \multicolumn{4}{c}{Opening offers} & \multicolumn{4}{c}{Later offers}\\
\cmidrule(lr){2-5}\cmidrule(lr){6-9}\cmidrule(l){10-13}
Would newly complete & $n$ & Acc. & Ctr. & Rej. & $n$ & Acc. & Ctr. & Rej. & $n$ & Acc. & Ctr. & Rej.\\
\midrule
Author only & 289 & 19 & 111 & 159 & 187 & 3 & 46 & 138 & 102 & 16 & 65 & 21\\
Recipient only & 140 & 24 & 78 & 38 & 51 & 2 & 18 & 31 & 89 & 22 & 60 & 7\\
Both & 81 & 10 & 35 & 36 & 48 & 3 & 11 & 34 & 33 & 7 & 24 & 2\\
Neither & 408 & 60 & 169 & 179 & 238 & 24 & 62 & 152 & 170 & 36 & 107 & 27\\
\midrule
All & 918 & 113 & 393 & 412 & 524 & 32 & 137 & 355 & 394 & 81 & 256 & 57\\
\bottomrule
\end{tabular}
\caption{Responses to all 918 acceptance-legal offers, by who would newly complete a color group and by position in the thread. Rows are offers, not independent trials. The same thread, pair and game contribute many rows, and prices, holdings and responders differ across rows. Re-tabulated from the retained offer ledger for this paper.}
\label{tab:cells}
\end{table}

Table~\ref{tab:cells} gives every response. Pooled over positions, an offer that completes only its author's group is rejected outright 159 times in 289 (55 percent) and accepted 19 times (7 percent). An offer that completes only the recipient's group is rejected 38 times in 140 (27 percent), countered 78 times and accepted 24 times (17 percent).

The pooled gap is partly a matter of position. Author-only offers are mostly a buyer's opening bid for the last deed it needs, and 187 of 289 open a thread. Recipient-only offers are mostly a seller naming its price in reply, and 89 of 140 follow a counter. Later offers are rejected far less often in every row. Within position the rejection gap between the two one-sided rows is about half the pooled one. Among opening offers it is 138 of 187 against 31 of 51 (74 against 61 percent). Among later offers it is 21 of 102 against 7 of 89 (21 against 8 percent).

Measured against offers that complete nothing, author-only offers are rejected somewhat more often at each position (74 against 64 percent of opening offers, 21 against 16 percent of later ones). The larger difference is in acceptance as written. Of 187 opening offers that would complete only the author's group, three are accepted, and two of the three are Grok tool calls whose text declined the offer. Of 238 opening offers that complete nothing, 24 are accepted. An opening offer that would complete only its author's group was taken as written once in this cohort by a response whose text and call agree.

The contrast is not uniform across systems (Table~\ref{tab:styles}). Sol, Gemini and Grok each reject author-only offers at a higher rate than recipient-only offers and accept them at a lower rate. Opus does neither. It rejects 3 of 52 against 6 of 42, accepts 7 of 52 against 5 of 42, and counters most offers of either kind (42 of 52 and 31 of 42). Opus answered 162 of its 189 offers through the action-only path, so this style cannot be separated from that condition. The direction of the rejection gap holds in all five seed blocks and in 17 of the 19 games that contain both kinds of offer. The acceptance gap holds in four blocks and 13 games.

\begin{table}[t]
\centering\small
\begin{tabular}{@{}lrrrrrrl@{}}
\toprule
 & \multicolumn{4}{c}{All 920 responses} & \multicolumn{2}{c}{Rejected} & \\
\cmidrule(lr){2-5}\cmidrule(lr){6-7}
Recipient & Received & Acc. & Ctr. & Rej. & Author only & Recipient only & Response condition\\
\midrule
Sol & 235 & 30 & 140 & 65 & 31 of 83 & 3 of 30 & standard\\
Gemini & 242 & 26 & 36 & 180 & 63 of 76 & 15 of 31 & standard\\
Opus & 189 & 29 & 141 & 19 & 3 of 52 & 6 of 42 & action-only in 162\\
Grok & 254 & 28 & 77 & 149 & 62 of 78 & 14 of 37 & text and call disagree in 11\\
\bottomrule
\end{tabular}
\caption{Responses by recipient. The first four columns use all 920 responses. The rejection columns use acceptance-legal offers in the named row of Table~\ref{tab:cells}. No response is an engine default. Seven of Grok's 28 acceptances come from tool calls whose text rejected or countered.}
\label{tab:styles}
\end{table}

The four systems accept between 26 and 30 offers each and differ in how they decline. Sol and Opus mostly counter. Gemini and Grok mostly reject, and they supply 125 of the 159 author-only rejections. Gemini rejects 62 of 67 opening author-only offers and also 83 of 95 opening offers that complete nothing, so for Gemini the row label adds little to the fact that the offer opens a thread.

\subsection{Counters reprice, with or without completion}
\label{sec:counters}

A recipient who counters an offer usually keeps the deed that was asked for and changes the price. Table~\ref{tab:counters} compares each counter with the offer it answers. Of 111 counters to author-only offers, 92 still hand over every deed the author asked for, and 79 are pure repricings that keep the whole bundle and raise the cash, by a median \$125. The same is true where no completion is involved. Of 110 counters to offers that asked for a deed and complete nothing, 87 keep every requested deed and 81 are pure repricings, at a median \$100.

\begin{table}[t]
\centering\small
\begin{tabular}{@{}lrrrr@{}}
\toprule
Offer answered & Counters to a request & Keep every & Pure repricing, & Median\\
would complete & for a deed & requested deed & cash raised & increase\\
\midrule
Author only & 111 & 92 & 79 & \$125\\
Neither & 110 & 87 & 81 & \$100\\
Both & 35 & 26 & 25 & \$150\\
Recipient only & 32 & 24 & 24 & \$112.50\\
\bottomrule
\end{tabular}
\caption{What a counter changes, for the 288 counters to an offer in which the author asked for at least one deed. A pure repricing keeps every deed and card on both sides and changes only the cash. The median is taken over pure repricings. Re-tabulated from the offer ledger for this paper.}
\label{tab:counters}
\end{table}

Counters are therefore how these systems bargain over price in general, and the share that keeps the deed on the table does not distinguish a deed that completes a rival's group (83 percent) from one that does not (79 percent). The size of the markup does not distinguish them either. Where a cash-for-deed offer is answered by a pure repricing, the counter asks a median 1.56 times the cash offered for a completing deed (47 offers) and 1.54 times for a non-completing one (56 offers).

The asking level differs more, and by system. In the same pure repricings Sol asks a median 2.9 times list price for a deed that would complete the buyer's group and 2.0 times for one that would not (22 and 27 counters). Grok asks 3.2 and 2.2 times (7 and 14). Opus asks 1.1 and 1.8 times (18 and 13), the reverse order. In offers of cash for one unmortgaged color deed, the buying author bids a median 1.59 times list when the deed completes its group (80 offers) and 1.41 times when it does not (118 offers). These are offered terms on different deeds at different moments. The accepted one-deed sales are too few to compare (between 4 and 12 per cell). The record shows higher asks and bids around completing deeds for some systems and no identified premium.

Cash direction adds nothing to this. The side that would complete a group is the side paying cash in 373 of the 429 one-sided offers, but the side that receives more deeds is the side paying in 309 of 313 offers that complete nothing. The paying side is the buying side throughout.

\subsection{The offers that were accepted}

Nineteen offers that would complete only the author's group were accepted. They are few enough to describe. Three are Grok tool calls that accepted while the text of the same response rejected, and all three completed a group for Opus. In eight the recipient held under \$100 when it accepted. Seven were accepted by Opus, all through the action-only path. Sixteen of the nineteen followed at least one counter.

Thread outcomes show the same thing from the proposer's side. Of 187 threads that open with an offer completing only the initiator's group, 26 settle. In 24 of the 26 the initiator still completes its group. It pays more cash than it first offered in 18 of the 24, and the median change is \$97.50 in the recipient's favor.

In this cohort, then, a proposal to complete one's own group is almost never accepted as first written. The threads that settle mostly settle on terms the recipient has improved, and the acceptances include eight by recipients short of cash and three by calls that their own text contradicts.

\subsection{Stated reasons}
\label{sec:reasons}

Seventy-four acceptance-legal offers that would have completed the recipient's group were rejected outright. Thirty-eight would have completed only the recipient's group and 36 would also have completed a group for the author. Grok wrote 31 of the rejections, Gemini 31, Opus 8 and Sol 4. All 74 have a public message and 66 a private note. One language-model reading finds a stated consequence for the rival in 63, such as the rival completing a set, gaining cash to build, or surviving. A second reading agrees on 61. Thirty-five of the first reader's 63 are among the 36 rejections where the rival's completion is a fact of the offer, and 28 are among the 38 where it is not. Thirty-four of the 74 come from two games.

We conclude that Grok and Gemini commonly state a reason about the rival when they refuse their own completion, and we put no rate on it. The reading does not check the arithmetic in a note or show that the stated reason produced the rejection, and two notes adopt the wrong role in the trade (G15 d106, G17 d186). For Sol and Opus the coded population is too small to say anything.

Low acceptance does not settle the worry, raised in studies of bilateral bargaining, that language models agree too readily~\cite{cosentino2026counterparty,shah2025rationalis}. A system could accept one offer in eight and still accept the wrong ones, and the cohort has no valuation against which to call accepted terms generous. The record shows only that agreement is not the default response of any of the four.

\subsection{Two bargaining incidents}

\IncidentTitle{g2}{A completion priced, refused, and later sold the other way}

At turn 45 of game 2, Opus holds two of the three yellow deeds and \$845. Sol holds the third, Marvin Gardens, and \$161. Opus already owns the light-blue group. It offers \$325 for Marvin Gardens. Sol counters:

\TranscriptQuote{g2-d124-price}

\TranscriptQuote{g2-d124-note}

The cash figures in the note are right. Paying \$325 would leave Opus \$520 and paying \$650 would leave \$195. The capacity claim is wrong. A yellow house costs \$150, so \$520 buys three houses, and nine would cost \$1,350. Opus counters at \$400 and Sol at \$600. Opus then writes:

\TranscriptQuote{g2-d127-final}

\TranscriptQuote{g2-d127-ceiling}

Sol answers:

\TranscriptQuote{g2-d128-keep}

At turn 69 Opus has \$17 and Sol \$435. Sol offers two railroads for Opus's two yellow deeds. Opus, now on the action-only path, counters, and Sol accepts:

\TranscriptQuote{g2-d197-reverse}

\TranscriptQuote{g2-d198-note}

\IncidentOutcome{The turn-45 thread ends without a trade. At turn 69 Sol receives the two yellow deeds for the railroads and \$100. Three turns later Opus is eliminated owing Gemini \$1,300, and the railroads pass to Gemini with the rest of its estate. Sol spends \$300 on yellow houses, collects \$44 on the group and none on a built deed, and recovers \$150 by selling the houses. At turn 85 it owes Gemini \$1,300 with \$77 and no asset left to sell or mortgage. Gemini inherits the yellow group, spends \$900 on it the next turn and wins at turn 102.}

\IncidentReading{Sol prices the deed by what it would let Opus build, and says so in both channels, though its count of houses is three times too high. The \$450 that Opus calls final is \$25 under the ceiling in its own note. The task permits that bluff, and here it ends the thread, so whether Sol would have sold at \$475 is unknown. The later sale runs the other way and happens when the seller has \$17. The sequence contains every step this paper studies. A completion is priced and refused, then sold by a holder short of cash, built on for almost no return, and passed through bankruptcy to the eventual winner. It does not show that either refusal or either sale was a mistake.}

\IncidentSource{G2 d124--128, turn 45, and d197--198, turn 69. Sol's responses and Opus's turn-45 responses are standard first attempts. Opus's d197 message is the selected second attempt of an action-only chain. Each quoted message appears in every other player's next delivered prompt. Cash, estates and spending are from the offer, flow and spending ledgers and the d215 and d239 liquidation menus.}

\IncidentTitle{swap}{Refusing a monopoly because the other side could build first}

In game 15, turn 63, Sol has \$730, two of the three green deeds and two of the three yellow deeds. Grok has \$333, the third green and the third yellow. Gemini has three houses on each orange property. Sol proposes a swap that would complete yellow for itself and green for Grok:

\TranscriptQuote{fx-g15-offer}

Sol's private note on the same action states the asymmetry it is counting on:

\TranscriptQuote{fx-g15-offer-note}

Grok rejects, and says why in public:

\TranscriptQuote{fx-g15-refusal}

Its private note gives the arithmetic:

\TranscriptQuote{fx-g15-refusal-note}

The cash figures match the ledger. Accepting would leave Sol with \$830 and Grok with \$233. A yellow house costs \$150 and a green house \$200, so \$830 buys five yellow houses and \$233 buys one green house but not one on each deed, which is what the note says. Sol drops the cash request and adds a third deed, is rejected again, and then offers to pay:

\TranscriptQuote{fx-g15-sweetener}

\TranscriptQuote{fx-g15-final-refusal}

\IncidentOutcome{No trade occurs. Grok is the first player eliminated, eighteen turns later, and its deeds, including the yellow and green blockers, pass to its creditor Gemini. Sol is eliminated last and never completes either group. Gemini wins.}

\IncidentReading{This offer would complete a group for both sides, and Grok declines its own completion three times, the last time with \$250 added. Each refusal gives a reason about what Sol could do with the other half of the swap. Sol's note shows it counting on the same asymmetry, and its third message answers Grok's stated objection with figures. Both sides reason about relative ability to build from public cash, which is accounting and not inference about hidden beliefs. The record does not show that refusing was the better choice. Grok went bankrupt first with the blockers in hand, and a third player inherited them.}

\IncidentSource{G15 d168--175, turn 63. All quoted responses are standard first attempts. Each quoted message appears in every other player's next delivered prompt. Cash and completion consequences are from the offer ledger, elimination order and estates from the flow ledger. Appendix~\ref{app:incidents} gives a second refusal of this kind from a two-player position (G11) and a ``final'' price that settled (G7).}

% G19 follows here unchanged from the Fable draft (Incident "abort"), with
% its reading sentence "cannot be read as a plan to mislead" replaced by:
% "so the false message does not establish a plan to mislead."

% ---------------------------------------------------------------------
% R5. DEVELOPMENT (replacement subsections)
% Base: Fable "Similar completions, different conversion" with its figure and
% table; Opus "Most end_turn records are not verified holds" with its table.
% Replace the two subsections below. Keep the G20 incident from any draft.
% ---------------------------------------------------------------------
\section{From completion to development (replacement subsections)}
\label{sec:development}

\subsection{Why twenty groups were not built on}
\label{sec:unbuilt}

Forty-three of the 63 trade-completed groups receive construction in the owner's first management window. The other twenty have four kinds of history, and the record supports a different statement about each.

\paragraph{Nine windows never listed construction.} The engine lists construction on a group only when no deed in it is mortgaged, the owner's cash covers one house and the bank has stock. It lists a redemption only when cash covers it. In eight of the nine a group deed was mortgaged when the window opened, and in the ninth the owner held \$158 against a \$200 house. The nine are not alike, and reading the delivered menus separates them (Appendix~\ref{app:unbuilt}).

Four are Grok's. In each, redeeming the mortgaged deed is on the menu, and Grok's note gives a reason to keep the cash. At G14 d207 it holds \$602 and its note opens \InlineQuote{rv-g14-hold} At G16 d234 it holds \$523 and its note says \InlineQuote{rv-g16-hold} In those two windows cash covers the redemption and one house. These four are stated holds with the enabling action available.

Five are Opus's, all on the action-only path and so without a note. Opus enters them with \$7, \$31, \$158, \$268 and \$125. In one it redeems the mortgaged deed and ends the turn with \$75 against a \$200 house. In three the redemption is itself unaffordable from cash. No text says what Opus intended. Three of its five public messages announce the completed group.

None of the nine is shown to be impossible to build on. We read menus and cash and did not search the sequences of mortgages and redemptions that might have reached a house within the window. The supported statement is that construction was never offered on the menus the owner saw, for a stated reason in four windows and with little cash in five.

\paragraph{Four ended in an engine default.} Three are Opus windows in which the first attempt was refused and the action-only retry submitted an invalid build plan. One is a Sol window with two invalid attempts. The engine applied \End{} in each.

\paragraph{Four contain a text and tool-call disagreement.} All are Grok's. The received call ends the turn while the text of the same response names a build or a financing step.

\paragraph{Three are end-of-turn choices with construction listed.} Two are Grok's and one is Gemini's. Every action in these windows is a valid selected response.

Restricting to the 42 groups that were buildable when the window opened gives 34 built and eight not, of which three are defaults, three are disagreements and two are stated holds. The acquisition itself adds one more condition. Three of Opus's fifteen trade completions were accepted by Grok tool calls whose text declined the offer, and all three are among Opus's eight unbuilt windows.

\subsection{The default path}
\label{sec:defaults}

Across all management decisions the engine applied a default \End{} 66 times, always after failed construction attempts: 61 for Opus, four for Sol and one for Gemini. Fifty-five fall in two games. The failures share a cause. The engine checks the even-building rule after each entry of a build plan, so a plan such as three houses on one deed followed by three on the next is rejected although the final layout is even.

The earlier audit tested every one of these plans at its exact historical state. Seventy attempts carry a well-formed plan and all seventy are invalid as submitted. Sixty-three execute once each entry is split into single purchases and ordered evenly, with the same cash, bank stock and purchases, and seven cost more than the player's cash. Every one of the 61 Opus decisions contains a plan that executes this way. In sixty of them the retry had been spent on the provider refusal, and in G2 d129 the invalid plan came first and the refusal second. The other systems meet the same rule with a correction available and usually end the decision with some construction applied (Table~\ref{tab:recovery}).

These defaults establish one thing exactly. At 61 decisions Opus requested a purchase that was reachable in that state, the request was invalid as written, and the turn ended with nothing built. They do not show that the purchases were good investments, and they do not measure rent that was lost. Three of Opus's eight unbuilt first windows are defaults. The other five are the low-cash windows above.

Of 663 end-of-turn actions applied on menus that listed construction, 66 are these defaults and 24 are Grok decisions whose text named another action. The remaining 573 show no detected conflict, which is weaker than 573 decisions to hold cash. Sol and Gemini return no assistant text, so a conflict of this kind could not be seen for them.

% ---------------------------------------------------------------------
% R6. LIQUIDITY (replacement paragraphs)
% Base: Sol "Liquidation depends on legally accessible finance" with the
% figure; Opus Table "four eliminations not settled by the bound" without
% its "Menu" narrative. Keep the G12 incident (Opus or Fable version).
% ---------------------------------------------------------------------
\section{Liquidity and elimination (replacement paragraphs)}
\label{sec:liquidity}

\paragraph{The three payable debts.} Three eliminations were payable from the debtor's own assets, and each has an executed payment sequence from the earlier audit. The sale action was on the liquidation menu in all three, and two of the three notes count it. Sol at G12 d190 owes \$500. Its note adds \$225 from selling three houses, \$130 from the mortgages then listed and \$53 in cash, reaches \$408 and concludes that no liquidation can pay. It leaves out the \$340 of mortgages on its three red deeds, which become available once the houses are gone. Gemini at G17 d234 owes \$550. Its note adds \$500 from selling five houses and \$30 from one mortgage and stops at \$530. It leaves out \$375 of mortgages on the two dark-blue deeds under those houses. Both omissions are the same second step, collateral that appears only after a sale. The third case is different. Grok at G12 d249 works through the staged sales in its note and concludes that selling is correct, and the received tool call declares bankruptcy (Incident~\ref{inc:bankruptcy}).

\paragraph{Conflicts in the other direction.} G12 d249 is the only one of Grok's 139 text and tool-call disagreements in which the call declares bankruptcy. In seven other liquidation decisions the disagreement runs the other way. The text names bankruptcy and the call mortgages or sells. All seven precede eliminations whose debt exceeds the own-asset bound, so none of them changed whether the debt could be paid.

% In the G12 reading, replace "every text channel argues against the action"
% (Opus) and "cannot be read as a decision to surrender" (Fable) with:
% "The assistant text and the private note argue for selling, and the public
% message that accompanies the call announces the bankruptcy. The record
% does not establish a decision to surrender."

% ---------------------------------------------------------------------
% R7. INHERITANCE (complete replacement of the section)
% ---------------------------------------------------------------------
\section{Estates and the bargaining around them}
\label{sec:inheritance}

\subsection{Estates follow rent}

A bankrupt player's deeds go to its creditor with their buildings. Fifty-nine of the sixty estates pass to a player and one, the last in game 13, to the bank. The final estate of a game arrives when play is over, so the population that can matter for play is the forty eliminations before the last.

In 35 of those forty the creditor is the eventual winner. In 16 games one player is the creditor in both, and that player wins. The eventual winner receives 202 of the 227 deeds and \$7,148 of the \$7,843 in cash that these estates carry. They complete 52 color groups while the game is still running, 45 of them for the eventual winner. For scale, trades complete 63 groups and move 207 deeds in all twenty games. A further 69 groups are completed by the game-ending estates, which we do not count as acquisitions anyone could use.

The groups that estates complete earn little. Rent paid on complete groups totals \$95,278 of the \$121,154 paid in all. Groups completed by trade collect \$57,378 of it, groups completed by purchase from the bank \$23,070 and by auction \$3,358. Groups completed by an in-play estate collect \$11,472, and 14 of the 52 collect anything. The median such group arrives 25 turns before the game ends.

The order of events is therefore rent first and estate second. Fifty-seven of the sixty debts are rent on a developed group, the owner of that group becomes the creditor, and the estate adds to a position that was already collecting. Nothing here measures what inheritance contributes to winning, because the creditor is usually the player who already owned the group that caused the bankruptcy. The finding concerns the rule. Under intact inheritance a player who cannot pay transfers a portfolio to one specific rival, and every other player has a stake in which rival that is.

\subsection{Bargaining over who inherits}

The systems act on that stake. In game 1, Opus has \$18 and offers its mortgaged Boardwalk to Grok, who holds Park Place. Grok accepts at \$100 with this note:

\TranscriptQuote{fx-g1-boardwalk-note}

The note gives Grok's reason for paying above its own ceiling, which is that Opus's bankruptcy would hand the deed to a third player. One figure in the excerpt is off for this engine, where redeeming Boardwalk costs \$221 and not the standard \$220. After the price and \$20 of transfer interest Grok holds \$219, short of either. Opus is eliminated four turns later and its estate goes to that third player, Gemini.

% The G13 and G18 incidents follow. Use the Opus or Fable card sequence for
% each and replace their outcome and reading blocks with the four below.

\subsection*{Replacement outcome and reading for the G13 incident}
\IncidentOutcome{Opus pays \$200 plus \$16 in mortgage-transfer interest and Grok holds \$235. In the retained exact-prefix check, accepting \$130 and then mortgaging Ventnor gives Grok \$295, enough for the rent it named. Grok does not land on Short Line on its next roll. Its cash stays between \$179 and \$309 for the next 39 turns. Opus completes the green group at auction seven turns after the purchase. At turn 146 Grok pays \$200 on another of Sol's railroads, and at turn 150 it lands on Short Line with \$49 and nothing left to mortgage. The yellow group passes to Sol. Sol is later eliminated by a \$100 tax, and Opus wins.}

\IncidentReading{The seller prices a deed by the buyer's interest in the seller's survival, and the buyer pays that price. The statement that \$130 leaves Grok \$35 short is true of cash and leaves out a way to pay that Grok's own note records, and the returned reasoning says the omission is intended. The \$70 between counter and settlement is observed. That the omission produced it is not. Grok had framed the stakes before the omission, Opus's text never says the cheaper route is impossible, Opus was buying a deed for a group it completed seven turns later, and Grok's note sets its own floor at \$165, which reads as a bargaining position as much as a claim of need. The opening probability is also wrong, since Short Line is a 4-in-36 roll from that square. Grok played 43 more turns before the transfer it had described took place. The record cannot say how many of those turns the sale bought, because the feared roll did not come and Grok kept a mortgage it could have used.}

\subsection*{Replacement reading for the G18 incident}
\IncidentReading{The record shows a repeated assurance, an explicit acceptance of it, and a withdrawal. The reason in Opus's returned explanation is the seller's insolvency. Paying \$250 would extend the game of a player it describes as essentially out. That is Opus's stated reason, not a finding that the purchase had lost its value, since the deed would also have completed Opus's green group. The record does not show that the offer was insincere when made or that Sol gave anything up in reliance on it. The game has no object for a standing offer and none of the messages set an expiry. In Incident~\ref{inc:insurance} a third party pays a seller who invokes its own bankruptcy. Here the likely creditor declines to.}

\subsection*{Closing paragraph of the section}
\paragraph{Blocking deeds outlive their holders' solvency.} Several deeds bought or kept to deny a rival reach that rival, or another, through an estate. In game 18 Grok buys a mortgaged Illinois Avenue for \$100, saying this keeps Sol off the red group. Eighteen turns later Grok is bankrupt to Sol, and Illinois completes Sol's reds. In game 15 Ventnor, the deed Grok refused to trade, passes to Gemini with the rest of its estate. In game 13 the yellow group that the insurance sale was meant to keep from Sol reaches Sol. In game 2 the group Sol paid to complete reaches Gemini. These are sequences, and none shows what the holder would have gained by acting differently. They do show why solvency is something these systems bargain about.

% ---------------------------------------------------------------------
% R8. DISCUSSION (replacement of the first subsection and the table)
% ---------------------------------------------------------------------
\section{Discussion}
\label{sec:discussion}

\subsection{What the account says about an economic agent}

\paragraph{A deal's value was set after the deal.} Two systems turned negotiated groups into houses almost every time and two did not, and the groups built at the first opportunity collected nearly all the rent. The reasons differ by system and by window. Grok's notes argue for holding cash, and its tool call sometimes ended a turn its text wanted to use. Opus arrived at its completions with little cash, and at 61 decisions it requested an affordable purchase in a form the validator rejected after its retry had been spent on a provider refusal. Scoring these games by wins would attribute all of this to the model, and scoring them by deals closed would miss it. An evaluation of an agent that holds assets should record three things separately for each intended action, whether it was reachable, whether its encoding was valid, and whether it was applied. Results from one execution setup should not be read as properties of the model inside it, a point made elsewhere with matched comparisons~\cite{zhang2026dece,hardy2026reliability}.

\paragraph{Bargaining was contested, and the pooled contrast is not a rival effect.} No system treats agreement as the default, and an opening offer that would complete only its author's group was accepted as written once in 187 attempts by a response whose text and call agree. Three systems reject such offers more often than offers for their own completion, and two of them say why in terms of the rival's position. The same data hold three cautions that any reading of such tables should carry. Position in the thread accounts for about half of the pooled gap. One system shows no gap. And counters reprice every kind of deed in the same way. We cannot say the systems refuse the right offers. Incident~\ref{inc:swap} shows a refusal that was coherent, checkable and followed by the refuser's bankruptcy.

\paragraph{Liquidity is what the rules let a player reach.} Fifty-six of sixty debts could not have been paid, so these systems almost never give up with a way to pay. The two that did left out the same step, the collateral behind buildings they were about to sell. Game 7 shows the opposite error available to an analyst, where \$1,310 of nominal assets could reach \$440. Both directions argue for measuring an agent's liquidity as the money a legal sequence can raise.

\paragraph{When default transfers assets, solvency becomes something to trade.} Under intact creditor inheritance a bankrupt player hands a portfolio to one rival. Estates in this cohort went to the player already collecting rent, and the groups they completed earned little. The systems nevertheless argued about who would inherit, priced a deed as insurance against it, and declined to prolong a debtor. The Vending-Bench Arena study reports that misaligned messages become more likely when an agent's inventory is low~\cite{li2026arena}. Our incidents are selected and carry no rate, but both records suggest that scarcity deserves separate attention when agents trade.

\subsection{Coverage of the broader questions}

\begin{table}[p]
\centering\footnotesize
\renewcommand{\arraystretch}{1.12}
\begin{tabularx}{\linewidth}{@{}L{2.5cm}YL{1.7cm}L{3.6cm}@{}}
\toprule
Question & What this cohort supports & Level & Not established\\
\midrule
1. Strategic reasoning & Offers completing the author's group are rarely accepted, once as an opening offer. Two systems state rival consequences when refusing their own completion (Incidents~\ref{inc:g2}, \ref{inc:swap}). Payable debts are abandoned in 3 of 60 eliminations. & Census and cases & That any choice raised the chance of winning. Regret or exploitability. A rival effect on refusal.\\
2. Opponent modeling & Notes and messages compute rivals' cash, collateral and building capacity, sometimes wrongly (Incident~\ref{inc:g2}). A proposer answers a stated objection. Some forecasts are checkable. & Cases & Accuracy over a defined set of forecasts. Inference about hidden beliefs.\\
3. Negotiation, cooperation, deception & All 920 responses with completion, position and counter content. Four declining styles at similar acceptance. Bluffed limits that settle and fail, an omitted financing route, a withdrawn offer, a corrected false premise. & Census for actions, cases for speech & Effectiveness against a valuation. Agreement bias or its absence. A deception rate. Persuasion.\\
4. Long-horizon planning, memory & The path from completion to building to rent for every traded group. Cash at the first management window. One target carried over 31 turns, one plan lost at the response layer. Memory conditions documented. & Census of trajectories, cases for plans & A rate of plan completion. Whether reserves were adequate. Any benefit of memory.\\
5. Adaptation & An encoding error recurs in 9 of 50 next submissions after the validator message was shown. Auction limits are revised. & Observed sequences & Learning from feedback. Adaptation to an opponent.\\
6. Economic decisions, risk & Financing feasibility at all 60 eliminations. Nominal against reachable finance. Stated reserve holds with the enabling action available. Offered prices relative to list. & Census & The value of a purchase before the fact. Reserve policy. Risk-adjusted return. A completion premium.\\
7. Collective behavior & Estates concentrate in the player collecting rent. One creditor takes both pre-final estates in 16 games. Blocking deeds pass through estates. One open appeal to deny a rival cash. & Census of flows, one case & How often coalitions form or last. Any effect of an appeal.\\
8. Validity, generalization & Conditions that change what records mean are stated and counted. Ledgers reconcile with the engine. Seat and seed dependence is visible in three blocks. & Internal consistency & Transfer to other opponents, rules or runs. Any ranking.\\
\midrule
A. Strategic talk and action & Text and applied action disagree in 139 Grok decisions, including one payable bankruptcy. Stated plans are carried out, lost at the response layer, or never tested. Requested builds were reachable and invalid as written. & Census for one system, cases & A general relation between explanation and action quality.\\
B. General and multi-agent capability & Not evaluated. Four configurations and no independent capability measure. & None & The relationship itself.\\
C. Skill, luck, exploitation & Seven games follow the seat that bought an early monopoly from the bank, under different systems. Four follow one system across seats. Specific execution failures are located. & Descriptive & A decomposition of winning into skill, chance and opponents' errors.\\
\bottomrule
\end{tabularx}
\caption{What the cohort supports for each motivating question. ``Census'' means a count over a complete population of recorded events. ``Cases'' means checked incidents with no rate. Levels describe the evidence, not the systems.}
\label{tab:coverage}
\end{table}

% ---------------------------------------------------------------------
% R9. RELATED WORK (replacement paragraph)
% Base: Sol "Related work and current context" for versions and wording.
% ---------------------------------------------------------------------
\section{Related work (replacement paragraph)}

\paragraph{Monopoly and financial board games.} AgentMonopoly studies four model families at two sizes with persona instructions and reports that larger models do not consistently win more games~\cite{wu2026agentmonopoly}. Its prompt encourages bluffing and prohibits public messages that contradict the action taken, so the text and action disagreements we count would be instruction violations there and are not here. The paper states that a player is eliminated when its cash becomes negative. Our statements about its creditor rule and its trade interface come from an earlier inspection of its repository~\cite{agentmonopoly2026repo} and not from the paper. FinBoardBench evaluates Cashflow, Acquire and Monopoly with one focal model against three fixed LLM opponents~\cite{hu2026finboardbench}. Its Monopoly rules give a player creditor the debtor's property but require the buildings to be sold to the bank, its trade manager allows counteroffers, and it feeds execution errors back to the model up to five times before applying a fallback that it records as system-forced. It reports that models keep buying property when short of cash. Two contrasts follow for our results. Intact buildings make an estate worth more here than under its rule. And the engine defaults we count arose under a budget of one correction, which the refusal path had already spent. Cattle Trade combines auctions, bluffing and bargaining and finds that strategic coherence, such as spending efficiency and resource discipline, correlates with rank more than any single skill, and that two scripted agents outperform most of its LLMs~\cite{muller2026cattletrade}.

% ---------------------------------------------------------------------
% R10. LIMITATIONS (complete replacement, consolidated)
% ---------------------------------------------------------------------
\section{Limitations}
\label{sec:limitations}

\paragraph{A small, dependent cohort.} Twenty games come from five seeds. Decisions within a game share a board, rotations within a block share early dice, and the four cyclic rotations cover four of the 24 possible seat orders. Every system faced the same three opponents, whose model names were visible. The counts are censuses of this cohort. They carry no confidence intervals and support no ranking, seat effect or forecast for other opponents.

\paragraph{Unequal conditions.} Opus played 95 percent of its decisions through an action-only recovery path whose trigger is unknown. Grok's text and tool calls disagree in 139 decisions for reasons the record cannot locate beyond the received response. Token use differs by a factor of thirty across systems. Twenty-nine percent of prompts contain a garbled line of trade history. None of these was assigned as a treatment, so comparisons across systems are comparisons across these conditions, and correcting them was outside this study.

\paragraph{One rule set.} Intact creditor inheritance, per-entry even-building checks and a bankruptcy option that is always on the menu shape the results. The engine overcharges some mortgage redemptions by one dollar. Replay shows that the record is consistent with the implementation. It does not certify the implementation against an independent statement of the rules.

\paragraph{Observation without intervention.} Offers differ in price, contents, responder and position within every row we tabulate, so no contrast estimates an effect. The stratifications we report were chosen after seeing the data. Splits by eventual winner condition on an outcome the same events produce. No trade, refusal, purchase or bankruptcy is shown to be the best available play, because no reference policy or continuation value was computed. Menu and cash facts about unbuilt groups are not feasibility results.

\paragraph{Text and selection.} Private notes are absent for most Opus decisions, returned reasoning is available unevenly, and no channel is privileged. Semantic codes were assigned by language models that share a model or developer with participants, with one partly unblinded second reading and no human labels. The incidents were selected by the same kind of assistant to show the mechanism each section measures, and five of the seven involve Grok, whose channels are the most observable. They support no rate.

\paragraph{Inherited checks.} The engine executions behind the payable and unpayable classifications, the reordered build plans and the replay were performed in an earlier audit and were not rerun for this paper. The unpayable classification rests on an arithmetic bound in 56 cases and on an exhaustive search of a one-action graph in one. The retained checker does not certify impossibility where a bundled building sale is possible. Appendix~\ref{app:provenance} separates what was recomputed for this paper from what was inherited.

% ---------------------------------------------------------------------
% R11. CONCLUSION (complete replacement)
% ---------------------------------------------------------------------
\section{Conclusion}

In twenty complete games, four configured LLM systems completed similar numbers of color groups by trade and realized very different amounts from them. Two built on nearly every group at once. One held cash for stated reasons and sometimes sent a tool call that contradicted its own text. One arrived short of cash and lost 61 affordable build requests to an invalid encoding with no retry left. Bargaining was contested. An opening offer to complete only one's own group was accepted as written once by a response whose text and call agree, counters repriced every kind of deed alike, and two systems said they were refusing because of what the other side would gain. Nearly every bankruptcy was unpayable when it came, and the estates went to the player already collecting rent.

For evaluation, the practical consequence is that a multiplayer economic game has at least three records to read, of deals, of what was built and of who inherited, and that they can disagree. A win count, a deal count or a striking transcript read alone would have misdescribed all four systems. In seven of these twenty games the seat that bought an early monopoly from the bank won under whichever system sat there.

% ---------------------------------------------------------------------
% R12. APPENDIX REPLACEMENTS
% ---------------------------------------------------------------------
\section{Appendix replacements}

\subsection{The nine windows that never listed construction}
\label{app:unbuilt}

\begin{table}[H]
\centering\footnotesize
\setlength{\tabcolsep}{4pt}
\begin{tabular}{@{}llrlllr@{}}
\toprule
Window & Owner, group & Cash & Mortgaged group deed & Redemption & Applied & Redeem and\\
 & & & & on menu & & one house\\
\midrule
G1 d298 & Grok, orange & \$229 & St.\ James, Tennessee & yes & end turn & \$300\\
G1 d314 & Grok, dark blue & \$369 & Boardwalk & yes & end turn & \$421\\
G14 d207 & Grok, dark blue & \$602 & Park Place & yes & end turn & \$393\\
G16 d234 & Grok, orange & \$523 & New York & yes & end turn & \$211\\
\addlinespace
G8 d268 & Opus, brown & \$7 & Baltic & no & end turn & \$83\\
G12 d236 & Opus, light blue & \$31 & Connecticut & no & end turn & \$116\\
G15 d57 & Opus, dark blue & \$158 & none & n/a & end turn & \$200\\
G16 d248--249 & Opus, dark blue & \$268 & Park Place & yes & redeem, end turn & \$393\\
G20 d297 & Opus, dark blue & \$125 & Boardwalk & no & end turn & \$421\\
\bottomrule
\end{tabular}
\caption{The nine trade-completed groups whose first management window never listed construction, read from the delivered decision records. Cash is at the first management decision after completion. The last column is the engine's redemption charge for the mortgaged group deeds plus one house at list cost. It is arithmetic for orientation. No sequence was executed, so the table does not show which windows could have reached a house. Grok's four applied notes each give a reason to keep cash. Opus's five decisions are action-only and have no note.}
\label{tab:unbuilt}
\end{table}

\subsection{Research roles}

The project's human author directs the work and is responsible for the benchmark, the runs and the research questions. The analysis and writing were carried out by language-model research assistants under that direction. They were GPT-6 Astra, GPT-6.1 Sol, Claude Opus 5.5 and Claude Fable 5.1 for the 29 October audit and review, and the same four with Claude Fable 5.2 in place of 5.1 for this rewrite, in which each assistant wrote a complete draft before reading the others. GPT-6.1 Sol and Claude Opus 5.5 are also two of the four evaluated models, and the Astra and Fable assistants share developers with them. No review pass was blinded to this. No semantic judgment in the paper is a human annotation. Mechanical counts are reproducible from the scripts released with the paper.

\subsection{What was recomputed for the integrated text}
\label{app:provenance}

\begin{table}[H]
\centering\footnotesize
\renewcommand{\arraystretch}{1.1}
\begin{tabularx}{\linewidth}{@{}YL{2.2cm}L{5.4cm}@{}}
\toprule
Quantity & Status & Source\\
\midrule
Table~\ref{tab:cells}: cells by response and position; block and game directions; Table~\ref{tab:styles} & Recomputed & Retained offer ledger, 920 rows\\
Table~\ref{tab:counters}: deed retention, pure repricing, normalized asks; one-deed bids relative to list & Recomputed, new & Offer ledger joined within threads; board list prices\\
Deed direction against cash direction (309 of 313; 373 of 429) & Recomputed, new & Offer ledger\\
The nineteen accepted author-only offers; thread outcomes & Recomputed & Offer ledger joined to the conflict list\\
Split of the 74 coded rejections by row and recipient & Recomputed from first-reader codes & Retained rejection-reason file. The 61 agreed codes are inherited\\
Table~\ref{tab:unbuilt}: menus, cash and notes in nine windows & Read for this paper & Canonical decision records\\
Liquidation menus and notes at the four eliminations above the bound; seven reverse disagreements & Read for this paper & Canonical decision records and the conflict list\\
Grok's cash from G13 d250 to d311 & Read for this paper & Delivered prompts\\
Routes, estates, rent by route, seed-block purchases & Recomputed & Episode ledger, typed flows, cohort table\\
New excerpts (G2, G11, G15, G1, G16) and delivery of every public one & Re-verified & Canonical decision records; each message found in every other player's next prompt\\
Payment sequences, reordered build plans, the G7 search, replay, state reconciliation & Inherited & 126-case certificate suite of 29 October; not rerun\\
Formatter, delivery and memory-window counts; token and latency figures & Inherited & Earlier audit\\
Rival-consequence agreement (61 of 74); 240-decision coding & Inherited & Two and one language-model readings\\
\bottomrule
\end{tabularx}
\caption{Status of the quantities in the integrated text.}
\end{table}

% ---------------------------------------------------------------------
% R13. MINIMAL INSERTIONS FOR integrated/ (snapshot of 11 November, 05:49 UTC,
% manuscript.pdf sha256 prefix 0feea7e4c8014626)
% If the parent keeps the current integrated sections instead of taking R4,
% R5 and R7 whole, these sentences close the findings that still apply to
% that text. Each item names the file and the sentence it replaces or
% follows. integration_proposal.md section 12 maps items to findings.
% ---------------------------------------------------------------------
\section{Minimal insertions for the current integrated text}

\paragraph{13a. Abstract, in place of the sentence on 111 counteroffers.} An opening offer that would complete only its author's color group is accepted 3 times in 187, twice by a tool call whose own text declined it. Counteroffers keep the requested deed and raise the cash at similar rates whether or not the deed completes a group.

\paragraph{13b. Bargaining, in place of the two sentences on 373 of 429.} Cash direction carries no information about completion. The completing side pays in 373 of the 429 one-sided offers, and in offers that complete nothing the side receiving more deeds pays in 309 of 313.

\paragraph{13c. Bargaining, after the first-offer and later-offer counts.} The sharper difference within position is acceptance as written. Three of 187 opening author-only offers are accepted, two of them by Grok calls whose content declined, against 24 of 238 opening offers that complete nothing.

\paragraph{13d. Bargaining, in place of ``Thus many recipients keep a rival's completion available \ldots{} not a completion-specific premium estimate.''} Counters to offers that complete nothing do the same. Of 110 that answer a request for a deed, 87 keep every requested deed and 81 are cash repricings of the unchanged bundle, at a median increase of \$100. Recipients therefore keep a requested deed on the table at about the same rate whether or not it completes the asker's group (83 and 79 percent). The counter content is evidence of bargaining over price in general. It is not evidence of a response to completion.

\paragraph{13e. Bargaining, after the sentence on sixty-one of seventy-four.} Grok and Gemini wrote 62 of the 74 rejections, and 35 of the first reader's 63 rival-consequence codes fall among the 36 rejections in which the offer would also have completed the author's group. The coded reasons therefore describe two systems, and about half of them restate a fact of the offer.

\paragraph{13f. Bargaining, after the seller-cash medians.} Opus is the seller in 21 of the 48 and in 13 of the 21 below \$100. Without Opus sellers the medians are \$287 and \$642, and 8 of 27 sellers of a completing deed are below \$100. The direction holds without Opus, and most of the low-cash sales are one system's.

\paragraph{13g. G2 incident, after ``the direction reverses''.} Opus now holds \$17 and Sol \$435. To the finding add: Opus is eliminated three turns after the sale, and the two railroads pass to Gemini with its estate.

\paragraph{13h. Development, in place of ``Seven of the nine no-option groups \ldots{} before redemption or building.''} Eight of the nine have a mortgaged group deed when the window opens, and in the ninth the owner holds \$158 against a \$200 house. Four are Grok's. In each, redeeming the deed is on the menu and the applied note argues for keeping cash, and in two of them cash covers the redemption and a house. Five are Opus's, on the action-only path and without a note, entered with \$7, \$31, \$158, \$268 and \$125. In one of these Opus redeems the deed and stops at \$75. These are menu and cash facts read from the delivered records, not financing proofs.

\paragraph{13i. Finance, after the payable cases.} Sol's note stops at \$408 against \$500 and leaves out \$340 of mortgages behind its houses. Gemini's stops at \$530 against \$550 and leaves out \$375. Seven other Grok liquidation responses disagree in the opposite direction, with text that names bankruptcy and a call that mortgages or sells. All seven precede eliminations whose debt exceeds the bound.

\paragraph{13j. G13 finding, before ``Forty-three turns after''.} Grok does not land on Short Line on its next roll. Its cash stays between \$179 and \$309 for 39 turns, and it pays \$200 on another of Sol's railroads four turns before the end. Opus completes the green group at auction seven turns after the purchase.

\paragraph{13k. Discussion, after the sentence on seat 2 and the dark blues.} In a third block the player in seat 1 completes the orange group by purchase on turn 47 in three of the four rotations, under three configurations, and wins all three. Seven games thus follow a seat and four follow a configuration.

\paragraph{13l. Related work, AgentMonopoly.} Cite the repository entry for the three implementation statements (single offers, elimination at negative cash, intact transfer). The paper states only that a player is eliminated when its cash becomes negative. For FinBoardBench add that its trade manager allows counteroffers and that it returns execution errors to the model up to five times before a recorded forced fallback.
